\documentclass[10pt,a4paper]{amsart}
\usepackage[margin=19mm]{geometry}
\usepackage{amsmath,amssymb,mathtools}
\usepackage[scr=rsfs,cal=euler]{mathalfa}
\usepackage{xfrac}
\usepackage[unicode,hidelinks,bookmarks=false]{hyperref}
\newcommand{\ie}{i.e.\ }
\newcommand{\cf}{cf.\ }
\newcommand{\resp}{respectively\ }
\newcommand{\Subsection}{\subsection}
\providecommand{\nobreakdash}{\nobreak\hbox{-}}
\setlength{\emergencystretch}{2em}
\multlinegap0pt
\usepackage{amssymb}
\newtheorem{lemma}{Lemma}
\title{Spectrality of the Dirac Operator with Complex-Valued Periodic Coefficients}
\author{O. A. Veliev}
\date{Source: arXiv:2603.02884v1 (CC BY 4.0)}
\begin{document}
\maketitle

\noindent\emph{Source and licence.} Spectrality of the Dirac Operator with Complex-Valued Periodic Coefficients. Version arXiv:2603.02884v1 (CC BY 4.0), distributed by arXiv under the Creative Commons Attribution 4.0 International licence, http://creativecommons.org/licenses/by/4.0/, as recorded in the arXiv licence metadata for this exact version and shown on the abstract page https://arxiv.org/abs/2603.02884v1. Full licence notice: \texttt{licenses/arxiv-2603.02884-CC-BY-4.0.txt}.\par\smallskip

\noindent\textit{Editorial scope.} Counted unit: the article's lemma lemma (source0.tex lines 727--740). The article's complete original proof of it is delivered with it (source0.tex lines 742--828, one \texttt{proof} environment of 87 lines). No mathematics is written, repaired or re-derived: the statement and the proof are the article's own lines, in the article's own notation, and no retained line is substituted or re-flowed.  No further source-local context is needed: every cross-reference of every retained line resolves inside the two delivered spans.  Nothing was omitted from any retained span, so the package carries no omission record.  No retained line contains a citation, so the wrapper carries no bibliography and no bibliography entry.  No retained line contains a figure environment, an image-inclusion command or drawing code, so no image is delivered and the package declares no asset dependency.  Editorial formatting only, in addition: an AMS \texttt{amsart} preamble that restates the article's own declarations and macros used by the retained lines, namely the environments \texttt{lemma} on separate counters, together with the standard packages those lines need.  The retained environments print in this document as Lemma 1 in order of appearance, whereas the article prints the same items with the numbering of its own full document.  The complete native source of the article as distributed in the arXiv e-print for this exact version is bundled unchanged as \texttt{sources/195643/source0.tex}.
\begin{lemma}
If $\left\vert \operatorname{Re}b\right\vert \geq1$ and
\begin{equation}
\int\limits_{0}^{\pi}\left\vert a_{1}(x)\pm\left(  a_{2}(x)+b\right)
i\right\vert ^{2}+\left\vert a_{3}(x)-b\pm\imath a_{4}(x)\right\vert
^{2}dx\leq\frac{4}{\pi}, \tag{27}%
\end{equation}
then the circles
\[
C(1,2k\pm t\pm ib)=\left\{  \lambda\in\mathbb{C}:\left\vert \lambda-(2k\pm
t\pm ib)\right\vert =1\right\}
\]
for all $k\in\mathbb{Z}$ lie in the resolvent set of $L_{t}(Q)$.
\end{lemma}

\begin{proof}
Assume that there exists $\lambda\in C(1,2k+t+ib)$ which is an eigenvalue of
$L_{t}(Q)$. Then, by (26), we have
\begin{equation}
\left\vert \left(  \Psi_{\lambda},\varphi_{n,j,t}\right)  \right\vert ^{2}%
\leq\frac{\left\Vert (Q-Q_{b})\varphi_{n,j,t}\right\Vert ^{2}}{\left\vert
\lambda-(2n+\left(  -1\right)  ^{j}(t+ib))\right\vert ^{2}}. \tag{28}%
\end{equation}
Using the definitions of $Q,$ $Q_{b}$ and $\varphi_{n,j,t}$ together with
condition (27), we obtain
\[
\left\Vert (Q-Q_{b})\varphi_{n,j,t}\right\Vert ^{2}\leq\frac{2}{\pi^{2}},
\]
for all $n,j$ and $t.$ Therefore, if
\begin{equation}
\sum_{n\in\mathbb{Z}}\frac{1}{\left\vert \lambda-(2n+t+ib)\right\vert ^{2}%
}+\sum_{n\in\mathbb{Z}}\frac{1}{\left\vert \lambda-(2n-t-ib)\right\vert ^{2}%
}<\frac{1}{2}\pi^{2}, \tag{29}%
\end{equation}
then we obtain the following contradiction:
\begin{equation}
1=\left\Vert \Psi_{\lambda}\right\Vert ^{2}=\sum_{j=1,2;n\in\mathbb{Z}%
}\left\vert \left(  \Psi_{\lambda},\varphi_{n,j,t}\right)  \right\vert ^{2}<1,
\tag{30}%
\end{equation}
since $\left\{  \varphi_{n,j,t}:j=1,2;\text{ }n\in\mathbb{Z}\right\}  $ is an
orthonormal basis.

Thus, it remains to prove (29). First, consider the first summation in (29).
By assumption $\lambda=2k+t+ib+x+iy,$ where $x\in\lbrack-1,1]$ and $y=\pm
\sqrt{1-x^{2}}.$ Then
\begin{equation}
\left\vert \lambda-(2n+t+ib)\right\vert ^{2}=\left(  2s+x\right)  ^{2}%
+1-x^{2}=4s^{2}+4sx+1, \tag{31}%
\end{equation}
where $s=k-n$. Therefore, we have
\[
\sum_{n\in\mathbb{Z}}\frac{1}{\left\vert \lambda-(2n+t+ib)\right\vert ^{2}%
}=\sum_{s\in\mathbb{Z}}\frac{1}{4s^{2}+4sx+1}=1+
\]%
\[
\sum_{s\in\mathbb{N}}\frac{1}{4s^{2}+4sx+1}+\sum_{s\in\mathbb{N}}\frac
{1}{4s^{2}-4sx+1}=1+\sum_{s\in\mathbb{N}}\frac{2(4s^{2}+1)}{(4s^{2}%
+1)^{2}-16s^{2}x^{2}}.
\]
Since the function $(4s^{2}+1)^{2}-16s^{2}x^{2}$ on the interval $[-1,1]$
attains its minimum value at the points $x=1$ or $x=-1,$ in the both cases,
using the well-known equality
\[
\sum_{n=1}^{\infty}\frac{1}{(2n-1)^{2}}=\frac{1}{8}\pi^{2},
\]
we obtain
\begin{equation}
\sum_{n\in\mathbb{Z}}\frac{1}{\left\vert \lambda-(2n+t+ib)\right\vert ^{2}%
}\leq1+\sum_{n=1}^{\infty}\frac{1}{(2n+1)^{2}}+\sum_{n=1}^{\infty}\frac
{1}{(2n-1)^{2}}=\frac{1}{4}\pi^{2}. \tag{32}%
\end{equation}


Now, consider the second summation in (29). Since%
\[
\lambda-\left(  2n-t-ib\right)  =2k+t+ib+x+iy-\left(  2n-t-ib\right)
=2(k-n)+2t+x+2ib+iy,
\]
where $b=\alpha+i\beta,$ there exist $m\in\mathbb{Z}$ and $v\in(-1,1)$ such
that%
\[
\lambda-\left(  2n-t-ib\right)  =2m+v+\imath(2\alpha+y).
\]
Here, $\left\vert 2\alpha+y\right\vert \geq\sqrt{1-v^{2}},$ since $\left\vert
\alpha\right\vert \geq1.$ Therefore,
\begin{equation}
\left\vert \lambda-(2n-t-ib)\right\vert ^{2}=\left(  2m+v\right)  ^{2}%
+1-v^{2}=4m^{2}+4mv+1. \tag{33}%
\end{equation}
Now, instead of (31), we use (33), where $s$ and $x\in\lbrack-1,1]$ are
replaced by $m$ and $v\in(-1,1).$ Repeating the proof of (32), we obtain
\begin{equation}
\sum_{n\in\mathbb{Z}}\frac{1}{\left\vert \lambda-(2n-t-ib)\right\vert ^{2}%
}<\frac{1}{4}\pi^{2}. \tag{34}%
\end{equation}
Thus, (29) follows from(32) and (34). Hence, if the point $\lambda$ from the
circle $C(1,2k+t+ib)$ is an eigenvalue of $L_{t}(Q),$ then we arrive at the
contradiction (30). This means that the circle lies in the resolvent set of
$L_{t}(Q).$ In the same way we prove that the circle $C(1,2k-t-ib)$ also lies
in the resolvent set of $L_{t}(Q).$ The lemma is proved.
\end{proof}

\end{document}
